\medskip
\noindent\underline{D. $n=2m>0$.}

由 XII, 3.7, 群
\[
H^n_{\{x_0\}}\bigl(X_s,R\psi_{\bar\eta}(\Lambda(m))\bigr)
\quad\text{及}\quad
R^n\psi_{\bar\eta}(\Lambda(m))_{x_0}
\]
各有一个在正负号意义下确定的自然生成元, 分别记作 $\delta$ 和 $\delta'$. 设 $\phi$ 为从前一群到后一群的自然映射. 可规范 $\delta$ 和 $\delta'$ 的符号, 使
\begin{equation}\tag{2.2.5.2}
(\delta',\delta)=1;
\end{equation}
于是
\begin{equation}\tag{2.2.5.3}
\phi(\delta)=(-1)^m\,2\,\delta'.
\end{equation}

设 $Z$ 为二次型 (2.2.1.1) 的 Clifford 代数偶数部分的中心, 它是 $k(\eta)$ 的可分二次扩张; 设 $\epsilon$ 为惯性群 $I$ 的相应特征标, 其阶为 $1$ 或 $2$. 对 $\sigma\in I$, 有
\begin{equation}\tag{2.2.5.4}
\sigma\delta=\epsilon(\sigma)\delta
\end{equation}
以及
\begin{equation}\tag{2.2.5.5}
\sigma\delta'=\epsilon(\sigma)\delta'.
\end{equation}
在以 $\mathbb Z/2k$ 为系数的上同调中, 公式 (2.2.5.2)、(2.2.5.3) 及 (2.2.5.4) 蕴含变差为
\begin{equation}\tag{2.2.5.6}
\operatorname{Var}(\sigma)(a)=
\frac{\epsilon(\sigma)-1}{2}\,(-1)^m\,(a,\delta)\,\delta.
\end{equation}

\medskip
\noindent\underline{E. $n=2m+1$---消灭圈.}

由 XII, 3.7, 群
\[
H^n_{\{x_0\}}\bigl(X_s,R\psi_{\bar\eta}(\Lambda(m))\bigr)
\quad\text{及}\quad
R^n\psi_{\bar\eta}(\Lambda(m+1))_{x_0}
\]
各有一个在正负号意义下确定的自然生成元, 分别记作 $\delta$ 和 $\delta'$. 采用 D 的记号, 可规范其符号, 使
\begin{equation}\tag{2.2.5.7}
(\delta',\delta)=1;
\end{equation}
于是
\begin{equation}\tag{2.2.5.8}
\phi(\delta)=0.
\end{equation}

\medskip
\noindent\underline{F. $n=2m+1$---变差.}

惯性群 $I$ 在 $X_\eta$ 的上同调上的作用平凡. 这由 XII, 3.7 以及 $I$ 在 $Y_\eta$ 的上同调上的作用平凡得出; 后者是由于 $Y$ 在 $S$ 上固有且光滑.

对惯性群中的 $\sigma$, 必有
\begin{equation}\tag{2.2.5.9}
D(\sigma)(\delta')=\lambda(\sigma)\delta,
\qquad \lambda(\sigma)\in\Lambda(1).
\end{equation}
可加性公式(XIII, 1.4.3.4)化为
\[
\lambda(\sigma\tau)=\lambda(\sigma)+\lambda(\tau).
\]
设
\[
\epsilon:I\longrightarrow\Lambda(1)=\mu_k
\]
为满足
\[
\sigma\bigl(\sqrt[k]{t}\bigr)
=\epsilon(\sigma)\sqrt[k]{t}
\]
的同态, 其中 $t$ 为一个一致化元. 由驯惯性群的结构, $\lambda$ 必为 $\epsilon$ 的倍数. 因此, 对适当的 $\lambda_X\in\Lambda$(依赖于 $X/S$), 公式 (2.2.5.9) 可改写为
\begin{equation}\tag{2.2.5.10}
\operatorname{Var}(\sigma)(a)
=\lambda_X\epsilon(\sigma)(a,\delta)\delta
\end{equation}
稍后将确定 $\lambda_X$.
\footnote{印本在 (2.2.5.6) 及 (2.2.5.10) 中省略了 $(a,\delta)$ 中的逗号; 此处是 C 中所定义的配对.}